\chapter{CORRECTNESS PROOFS AND VALIDATION}\label{chap:correctness}

This chapter establishes the formal logical foundation of the proposed algorithms. We distinguish between \textit{Theoretical Guarantees}—mathematical certainties derived from the algorithm's structure—and \textit{Empirical Observations}—performance metrics that characterize the algorithm's behavior in practice.

\section{Correctness of the Log-Probability Transformation}

\begin{theorem}[Order-Preserving Log Transform]\label{thm:log_transform}
The path $P^*$ that minimizes $\sum_{(u,v)\in P} -\log p(u,v)$ is identical to the path that maximizes $\prod_{(u,v)\in P} p(u,v)$.
\end{theorem}

\begin{proof}
The negative logarithm $f(x)=-\log(x)$ is a strictly monotone decreasing function on the interval $(0,1]$. For any two paths $P_1, P_2$, the transformation preserves the relative ordering of path costs such that the path with the highest cumulative product of probabilities maps to the path with the lowest cumulative sum of negative log-probabilities~\citep{manning1999nlp}. Since $p(u,v) \in (0,1]$, all edge weights $w(u,v) \geq 0$, satisfying the non-negativity constraint required for Dijkstra's algorithm.
\end{proof}

\section{Correctness of Baseline Dijkstra}

\begin{theorem}[Baseline Optimality]\label{thm:baseline_optimal}
Upon termination, $\mathit{dist}[t]$ equals the true shortest-path distance $\delta(s,t)$ in the log-transformed graph.
\end{theorem}

\begin{proof}
We rely on the standard inductive proof for Dijkstra’s algorithm~\citep{cormen2009}. The algorithm maintains the invariant that for every node $u$ extracted from the priority queue, $\mathit{dist}[u] = \delta(s,u)$. Because all edge weights $w(u,v) \geq 0$, the extraction of a node $u$ with the minimum tentative distance guarantees that no shorter path can exist through nodes currently in the queue or yet to be discovered. 
\end{proof}

\section{Correctness of Probability-Pruned Dijkstra}

The Pruned Dijkstra variant introduces a threshold $\tau$, where $T = -\log \tau$. 

\subsection{Convergence and Conditional Optimality}

\begin{theorem}[Convergence]\label{thm:convergence}
As $\tau \to 0$, Algorithm 2 becomes identical to Algorithm 1.
\end{theorem}

\begin{proof}
Setting $\tau = 0$ results in a threshold $T = \infty$. The pruning condition $\mathit{new\_cost} > T$ is never met for any path with finite cost, effectively disabling the pruning mechanism and reverting the search to the baseline Dijkstra logic.
\end{proof}

\begin{theorem}[Conditional Optimality]\label{thm:conditional}
If the optimal path $P^*$ has a total probability $\Pi^* \geq \tau$, Algorithm 2 returns $P^*$ with an optimality gap $\Delta = 0\%$.
\end{theorem}

\subsection{The All-or-Nothing Property}

Unlike traditional approximation algorithms that may return a sub-optimal path within an $\epsilon$-bound, the proposed pruning method prioritizes \textbf{Optimality over Reachability}.

\begin{theorem}[All-or-Nothing]\label{thm:all_or_nothing}
Algorithm 2 either returns the globally optimal path or reports no path found. It never returns a sub-optimal path.
\end{theorem}

\begin{proof}
Since the pruning check only skips edges and never modifies the weights of edges it explores, the priority queue still processes nodes in non-decreasing order of cost. If the optimal path $P^*$ is not pruned, it will be found first. If any edge in $P^*$ is pruned because its prefix cost exceeds $T$, the target $t$ becomes unreachable via $P^*$. Consequently, the algorithm either finds the true optimum or terminates with no result; an optimality gap $\Delta > 0$ is impossible by construction.
\end{proof}

\section{Empirical Validation and Reachability}

Because we prioritize optimality, the primary "cost" of pruning is the potential failure to find a path. We define \textbf{Reachability ($R$)} as the core empirical metric to evaluate this trade-off.

\begin{definition}[Reachability]
The Reachability $R$ is the ratio of successful path-finding by the pruned algorithm compared to the baseline:
\\[-1em]
\[ R = \frac{\text{Count}(\text{Path Found}_{\text{Pruned}})}{\text{Count}(\text{Path Found}_{\text{Baseline}})} \]
\end{definition}



\subsection{Summary of Results}

The properties were verified across 3,000 experimental configurations. As shown in Table~\ref{table:correctness_summary}, the optimality gap remained $0.0000\%$ in all cases where a path was returned, while Reachability varied as a function of $\tau$.

\begin{table}[htbp]
\centering
\renewcommand{\arraystretch}{1.3}
\begin{tabular}{@{}lll@{}}
\toprule
\textbf{Property} & \textbf{Type} & \textbf{Validation Result} \\
\midrule
Order-Preserving & Theoretical & Verified ($<10^{-12}$ precision error) \\
All-or-Nothing & Theoretical & $\Delta = 0.0000\%$ in 100\% of found paths \\
Convergence & Theoretical & Identical output at $\tau = 0$ \\
Reachability ($R$) & \textbf{Empirical} & $R$ decreases monotonically as $\tau$ increases \\
Speedup & \textbf{Empirical} & Up to $18{,}882\times$ on real-world data \\
\bottomrule
\end{tabular}
\caption{Summary of Correctness and Performance Validation.}
\label{table:correctness_summary}
\end{table}

\section{Separation of Claims and Limitations}

\subsection{Proven Theoretical Claims}
The following are formally guaranteed by Theorems~\ref{thm:log_transform} through~\ref{thm:all_or_nothing}:
\begin{itemize}
    \item The algorithm is optimality-preserving (it never returns a "wrong" path).
    \item The log-transformation is mathematically sound for probability maximization.
\end{itemize}

\subsection{Empirical Observations (Expected Outcomes)}
The following depend on graph topology and data distributions and are not formally guaranteed:
\begin{itemize}
    \item \textbf{Search-Space Reduction:} The actual number of edges relaxed depends on the specific $\tau$ and the concentration of high-probability paths.
    \item \textbf{Optimal Threshold $\tau_c$:} Identifying a threshold that maximizes speedup while maintaining $R \approx 1$ is an empirical task performed via a sweep (Chapter 5).
\end{itemize}

\subsection{Limitations}
\begin{itemize}
    \item \textbf{No Reachability Guarantee:} For $\tau > 0$, there is no formal lower bound on $R$. In the worst case, a graph could have its only path pruned if $\Pi^* < \tau$.
    \item \textbf{Domain Specificity:} Results are currently validated on e-commerce clickstreams; performance on non-sparse or non-behavioral graphs is unknown.
\end{itemize}